\section{Completion and exact order}\label{sec:completion}

Equation \eqref{eq:sign-bridge} and \eqref{eq:central-identity} imply
\[
 \partial_\xi\Psi(\xi;m,r)<0\qquad(\xi<m<r).
\]
Proposition~\ref{prop:equivalence} therefore proves
Theorem~\ref{thm:main}.

We now prove directly that the classification stops at order four. Put
\[
 H_k=\det[\Phi^{(i+j)}(0)]_{i,j=0}^{k-1},\qquad
 s_j=\Phi^{(j)}(0).
\]
Since \(\Phi\) is even, simultaneous parity permutations of the rows and
columns give
\[
 H_4=AB,\qquad H_5=CB,
\]
where
\[
 A=s_0s_4-s_2^2,\qquad B=s_2s_6-s_4^2,
\]
and
\[
 C=\det\begin{pmatrix}
 s_0&s_2&s_4\\
 s_2&s_4&s_6\\
 s_4&s_6&s_8
 \end{pmatrix}.
\]

\begin{lemma}\label{lem:origin-pf5-sign}
For the Riemann kernel,
\[
 445<A<446,\qquad 189223<B<189228,
 \qquad -674000000<C<-614000000.
\]
In particular, \(H_4>0\) and \(H_5<0\).
\end{lemma}

\begin{proof}
For \(x=\pi n^2e^{2u}\), the \(n\)-th summand of \(\Phi(u)\) is
\(e^{u/2-x}P_0(x)\), where \(P_0(x)=2x(2x-3)\). If its \(j\)-th derivative
is \(e^{u/2-x}P_j(x)\), then
\[
 P_{j+1}=2xP_j'+(1/2-2x)P_j.
\]
Thus
\[
 s_j=\sum_{n\ge1}e^{-\pi n^2}P_j(\pi n^2).
\]
The first three modes are enclosed by rational alternating-series bounds.
Machin's identity
\(\pi=16\arctan(1/5)-4\arctan(1/239)\) supplies rational bounds for \(\pi\),
and argument division followed by the degree-29 and degree-30 alternating
Taylor sums supplies rational bounds for each exponential.

For the omitted modes, let \(S_j\) be the sum of the absolute coefficients of
\(P_j\). Since \(\deg P_j=j+2\), \(3<\pi<22/7\), and
\(e^3>\sum_{r=0}^8 3^r/r!>20\),
\[
 \sum_{n\ge4}|P_j(\pi n^2)|e^{-\pi n^2}
 \leq
 \frac{S_j(22/7)^{j+2}4^{2j+4}20^{-16}}
 {1-(5/4)^{20}20^{-9}}
 \qquad(j=0,2,4,6,8).
\]
Exact rational interval arithmetic then gives the stated bounds. The full
endpoints and the standard-library verifier are included in the
reproducibility archive.
\end{proof}

\begin{theorem}\label{thm:not-pf5}
The Riemann kernel is not \(\PF_5\).
\end{theorem}

\begin{proof}
For a \(C^{2k-2}\) function \(f\), repeated row and column divided differences
give, for increasing fixed tuples \(a,b\),
\[
 \lim_{h\downarrow0}
 \frac{\det[f(z+h(a_i-b_j))]_{i,j=0}^{k-1}}
 {h^{k(k-1)}V(a)V(b)}
 =\frac{(-1)^{k(k-1)/2}}
 {\prod_{r=0}^{k-1}(r!)^2}
 \det[f^{(i+j)}(z)]_{i,j=0}^{k-1}.
\]
Take \(k=5\), \(z=0\), and \(a_i=b_i=i\). The sign factor is one, both
Vandermonde factors are positive, and Lemma~\ref{lem:origin-pf5-sign} makes
the limit strictly negative. Hence
\(\det[\Phi((i-j)h)]_{i,j=0}^4<0\) for every sufficiently small positive
\(h\). These are translation minors at strictly increasing nodes.
\end{proof}

Theorem~\ref{thm:main} and Theorem~\ref{thm:not-pf5} prove
Corollary~\ref{cor:exact-order}.

The proof uses no multidimensional interval cover of the nonconfluent
configuration space. Earlier positive-tail Peano densities, collision cones,
Hermite boxes, and escape charts remain independent certified results, but
they are not premises here. Their role was superseded when the transport
numerator was recognized as the difference of expectations in
\eqref{eq:transport-expectation}.

\begin{remark}[Riemann Hypothesis boundary]
The theorem determines the kernel's exact finite Polya-frequency order and
makes no claim about the zeros of the Riemann \(\Xi\)-function. In particular,
the failure of \(\PF_5\) neither proves nor disproves the Riemann Hypothesis.
The next section proves, by an explicit strict-\(\PF_4\) separator, that no
Fourier real-rootedness conclusion follows from membership through order four
alone.
\end{remark}
